




\begin{thm}[Dekomposisi Polar]\label{0022287} Jika $T$ suatu operator pada ruang
 \index{dekomposisi polar}%
 \index{dekomposisi!polar}%
Hilbert $H$, maka terdapat isometri parsial $V$ pada $H$ sedemikian sehingga
 \begin{enumerate}
    \item[(i)] ruang awal $V$ adalah $(\ker T)^\perp$,
    \item[(ii)] ruang akhir $V$ adalah $\clo{\ran T}$, dan
    \item[(iii)] $T = V\abs T$.
 \end{enumerate}
Dekomposisi $T$ ini unik dalam arti bahwa jika $T = V_0P$ dengan $V_0$ suatu
isometri parsial dan $P$ suatu operator positif sedemikian sehingga $\ker V_0 = \ker P$, maka $P =
\abs T$ dan $V_0 = V$.
\end{thm}

\begin{proof} Lihat \cite{Conway:1990}, VIII.3.11;\allowbreak \cite{Davidson:1996}, teorema I.8.1;
\allowbreak \cite{KadisonR:1983}, teorema 6.1.2;\allowbreak \cite{Murphy:1990}, teorema 2.3.4;
atau~\cite{Pedersen:1995}, teorema 3.2.17. \ns
\end{proof}

\begin{cor} Jika $T$ adalah suatu
  \index{dekomposisi polar!dari operator invertibel}%
 \index{operator!invertibel!dekomposisi polar dari}%
 \index{invertibel!operator!dekomposisi polar dari}%
operator invertibel pada suatu ruang Hilbert, maka isometri parsial dalam dekomposisi polar
(lihat~\ref{0022287}) bersifat uniter.
\end{cor}

\begin{prop} Jika $T$ adalah suatu
 \index{dekomposisi polar!dari operator normal}%
 \index{operator!normal!dekomposisi polar dari}%
 \index{normal!operator!dekomposisi polar dari}%
operator normal pada suatu ruang Hilbert $H$, maka terdapat operator uniter $U$ pada $H$
yang berkomutasi dengan $T$ dan memenuhi $T = U\abs T$.
\end{prop}

\begin{proof} Lihat \cite{Pedersen:1995}, proposisi 3.2.20.\ns \end{proof}

\begin{exer} Misalkan $(S, \fml A,\mu)$ suatu ruang ukur $\sigma$-hingga dan $\phi \in
L_\infty(\mu)$. Tentukan dekomposisi polar dari
 \index{operator perkalian!dekomposisi polar dari}%
 \index{operator!perkalian!dekomposisi polar dari}%
 \index{dekomposisi polar!dari operator perkalian}%
operator perkalian $M_\phi$ (lihat contoh~\ref{C063527}).
\end{exer}









\section{Operator Kelas Jejak}

Mari kita mulai dengan meninjau beberapa sifat standar \emph{jejak} suatu matriks $n \times n$.

\begin{defn}\label{tr0113} Misalkan $a = [a_{ij}] \in \mathbf M_n$. Definisikan $\tr a$, yaitu
 \index{jejak!dari suatu matriks}%
 \index{matriks!jejak dari suatu}%
 \index{tr@$\tr$ (jejak)}%
\df{jejak} dari $a$, dengan
  \[ \tr a = \sum_{i = 1}^n a_{ii}. \]
\end{defn}

\begin{prop}\label{tr0117} Fungsi $\tr\colon \mathbf M_n \sto \C$ bersifat linear.
\end{prop}

\begin{prop}\label{tr0121} Jika $a$, $b \in \mathbf M_n$, maka $\tr (ab) = \tr (ba)$.
\end{prop}

Jejak bersifat invarian terhadap keserupaan. Dua matriks $n \times n$, yaitu $a$ dan $b$, disebut
 \index{serupa!matriks}%
 \index{matriks!serupa}%
\df{serupa} jika terdapat matriks invertibel $s$ sedemikian sehingga $b = sas^{-1}$.

\begin{prop}\label{tr0124} Matriks-matriks dalam $\mathbf M_n$ yang serupa mempunyai jejak yang sama.
\end{prop}

\begin{defn}\label{tr0211} Misalkan $H$ suatu ruang Hilbert separabel, $(e_n)$ suatu basis ortonormal bagi~$H$, dan $T$
suatu operator positif pada~$H$. Definisikan
 \index{jejak}%
 \index{tr@$\tr$ (jejak)}%
\df{jejak} dari~$T$~dengan
    \[ \tr T:=\sum_{k=1}^\infty \langle Te_k,e_k \rangle. \]
(Jejak tidak harus berhingga.)
\end{defn}

\begin{prop} Jika $S$ dan $T$ merupakan operator positif pada suatu ruang Hilbert separabel dan $\alpha \ge 0$, maka
   \[  \tr(S + T) = \tr S + \tr T \]
dan
   \[ \tr(\alpha T) = \alpha \tr T. \]
\end{prop}

\begin{prop} Jika $T$ suatu operator pada ruang Hilbert separabel, maka $\tr(T^*T) = \tr(TT^*)$.
\end{prop}

Untuk menunjukkan bahwa definisi jejak operator positif yang diberikan di atas tidak bergantung pada basis ortonormal yang dipilih bagi
ruang Hilbert tersebut, kita memerlukan suatu hasil (yang diberikan berikut ini) yang bergantung pada fakta yang muncul kemudian dalam catatan ini: setiap operator positif pada ruang Hilbert
memiliki akar kuadrat (positif) (lihat contoh~\futurexref{11.5.7}{X_sqroot_op}).

\begin{prop}\label{tr0221} Pada keluarga operator positif di suatu ruang Hilbert separabel $H$, jejak bersifat invarian terhadap ekuivalensi uniter;
yakni, jika $T$ operator positif dan $U$ uniter, maka $\tr T = \tr(UTU^*)$.
\end{prop}

\begin{proof}[\emph{Petunjuk untuk bukti}] Dalam proposisi sebelumnya, ambil $S = UT^{\frac12}$.    \ns
\end{proof}

\begin{prop}\label{tr0224} Definisi \emph{jejak}~\ref{tr0211} dari suatu operator positif pada ruang Hilbert separabel tidak bergantung
pada pemilihan basis ortonormal bagi ruang tersebut.
\end{prop}

\begin{proof}[\emph{Petunjuk untuk bukti}] Misalkan $(e^k)$ dan $(f^k)$ basis-basis ortonormal bagi ruang tersebut dan $T$ suatu operator positif
pada ruang itu. Ambil $U$ sebagai operator uniter unik yang memenuhi $e^k = Uf^k$ untuk setiap $k \in \N$.  \ns
\end{proof}
